\cotefile{128}
\scanpage{2}
\textcircled{1} \uline{Réductions via le conducteur}

\uline{Proposition} 1 \quad $A$ un anneau, $I$ un idéal de $A$, $B$ une $A$-algèbre \deleted{contenant $A$,} telle que l'application $\lambda\mapsto\lambda 1_B$ de $I$ dans $B$ soit \uline{injective}, et que son image soit \deleted{\unsure{telle que}} un idéal de $B$. \inserted{[nécessairement égal à $IB$]}

Pour tout entier $n\geq 0$ considérons l'anneau \inserted{$\unsure{C|\mathbb C}^{n}=\overset{n+1}{\bigotimes}_{A}B$} et l'application $\lambda\mapsto\lambda\deleted{K}1_{C^{n}}$ de $I$ dans $C^{n}$ ; \deleted{Alors} cette application est injective, et son image est un idéal. \deleted{De façon précise, soit $A_0=A/I$, $B_0=B\otimes_A A_0=B/IB=B/I1_B$, alors l'application $I\longrightarrow K^n$ précédente, et l'application $K_0^n=\overset n\otimes_{A_0}B_0\longleftarrow K^n=\overset n\otimes_A B$ donnent lieu à une suite exacte}

En effet, on a
\[
\begin{aligned}
\lambda(b_{\unsure{1|0}}\otimes\cdots\otimes b_n)&=(\lambda b_0)\otimes b_1\cdots\otimes b_n\\
&=(\lambda_{\unsure{1|0}}\deleted{\unsure{\otimes|\times}}1_B)\otimes b_2\cdots\otimes b_n
  \qquad\text{pour }\lambda_0\in I\ \text{conv.}\\
&=1_B\otimes[\lambda_1(b_2\otimes\cdots\otimes b_n)]\deleted{\unsure{=}}\\
&\hspace{-4em}\text{d'où le résultat par récurrence sur }n\\
&\hspace{-4em}\text{(trivial pour }n=0\text{)}\\[-.6ex]
&=1_B\otimes\mu(\overbrace{1_B\otimes\cdots\otimes1_B}^{n})
  \qquad\mu\in I\ \text{conv.}\\
&=\mu\deleted{\unsure{\otimes}}(\underbrace{1_B\otimes\cdots\otimes1_B}_{n+1})
\end{aligned}
\]
Pour l'injectivité, on compose
\[
I\mathrel{\unsure{-|\longrightarrow}}\overset n{\bigotimes}_A B\longrightarrow B,
\]
on trouve une application inj. par hyp., elle \unsure{l'étant|l'était} \unsure{avant}.
% archive sheet number: 1
\scanpage{3}
\deleted{\uline{Lemme} 1 \quad $A$ anneau, $I$ \deleted{\unsure{est|et}} un idéal de $A$, $B$ et $C$ deux $A$-algèbres, \deleted{telles que $B/\operatorname{Im} A$ et $C/\operatorname{Im} A$ sont annulés par $I$} \inserted{on \unsure{suppose}} $A\longrightarrow B$ injectif et $IB=I$, de $=$ $A\longrightarrow C$ injectif et $IC=I$. Alors $A\longrightarrow B\otimes_A C=D$ est injectif et}

\uline{Corollaire} 1 \quad Comme $I.1_{\unsure{C|\mathbb C}^{n}}=IC^n$, et $C^n/IC^n=C_0^n$ $(=\overset{n}{\bigotimes}_{A_0}B_0$, où $A_0=A/I$, $B_0=B/IB)$ on trouve On a une suite exacte (pour tt ens fini d'indices)
\[
0\mathrel{\unsure{-|\longrightarrow}} I\longrightarrow\overset E{\bigotimes}_{A}B\longrightarrow\overset E{\bigotimes}_{A_0}B_0\longrightarrow0
\]
Si on a une application d'ens finis \inserted{$E\mathrel{\mathop{\unsure{-|\longrightarrow}}\limits^\sigma}F$}, on a commutativité dans le diagramme \unsure{correspondant|correspondants}
\diagram{\begin{tikzcd}[ampersand replacement=\&,column sep=small]
0\arrow[r,no head]\& I\arrow[r]\arrow[d,equal,"\mathrm{id}"] \& \overset E{\bigotimes}_A B\arrow[r]\arrow[d,"\sigma_{B/A}"] \& \overset E{\bigotimes}_{A_0}B_0\arrow[r]\arrow[d,"\sigma_{B_0/A_0}"] \& 0\\
0\arrow[r,no head]\& I\arrow[r]\& \overset F{\bigotimes}_A B\arrow[r]\& \overset{\unsure{E|F}}{\bigotimes}_{A_0}B_0\arrow[r,no head]\& 0
\end{tikzcd}}
où les flèches verticales sont les hom simpliciaux déduits de $\sigma$, d'où une suite exacte de \deleted{groupes} \inserted{$A$-modules} simpliciaux
\[
0\longrightarrow C_I\longrightarrow C(B/A)\deleted{\unsure{C(B/A)}}\longrightarrow C(B_0/A_0)\longrightarrow0
\]
où $C_I$ est le $A$-module simplicial \uline{trivial} défini par $I$ (f\deleted{\unsure{cteur \uline{constant}|oncteur \uline{constant}}} $E\mapsto I$).
% archive sheet number: 2
\scanpage{4}
\uline{Corollaire} 2 \quad L'hom $K(B/A)\longrightarrow K(B_0/A_0)$ définit un isomorphisme
\[
H^i(B/A)\xrightarrow{\sim}H^i(B_0/A_0)\qquad\text{pour }i\ne0
\]
et une suite exacte
\[
0\longrightarrow I\longrightarrow H^0(B/A)\longrightarrow H^0(B_0/A_0)\longrightarrow0
\]
% margin note written sideways
\margin{d'où $\operatorname{Coker}(A\longrightarrow H^0(B/A))\simeq\operatorname{Coker}(A_0\longrightarrow H^0(B_0/A_0))$}

\inserted{\uline{Cor. 3} a) Pour que $H^i(B/A)=0$ il f. et s. que $H^i(B_0/A_0)=0$ ($i\ne0$ donné)

\deleted{\unsure{b.}} b) $A\longrightarrow\deleted{\unsure{B}}H^0(B/A)$ inj (surj) sss $A_0\longrightarrow H^0(B_0/A_0)$ inj (surj)}

Plus généralement, tensorisant par un $A$-module $M$ plat, on trouve

\uline{Corollaire} \deleted{\unsure{3}}4 \deleted{\unsure{bis}} \quad Soit $M$ un $A$-module \uuline{plat}. On a une suite exacte de \deleted{complexes} \inserted{$A$-modules simpliciaux}
\[
0\longrightarrow C_{M\otimes_A I}\longrightarrow K(B/A,M)\longrightarrow K(B_0/A_0,M_0)\longrightarrow0
\]
d'où des isomorphismes
\[
H^i(B/A,M)\simeq H^i(B_0/A_0,M_0)\qquad i\ne0
\]
et une suite exacte
\[
0\longrightarrow \unsure{\underline M|M}\otimes_A I\longrightarrow H^0(B/A,M)\longrightarrow H^0(B_0/A_0,M_0)\longrightarrow0
\]
% archive sheet number: 3
\scanpage{5}
\uline{Corollaire 5} \quad Supposons $B_0$ fid. plat sur $A_0$. Alors pour tt $A$-module plat \inserted{$M$}, on a
\[
\left\{\begin{aligned}H^i(B/A,M)&=0 &&i\ne0\\
H^0(B/A,M)&\xleftarrow{\sim}M\end{aligned}\right.
\]

\uline{NB} \quad Ce résultat est \uline{équivalent} au résultat correspondant pour $M=A$, i.e. à
\[
\left\{\begin{aligned}H^i(B/A)&=0 &&i\ne0\\
H^0(B,A)&\xleftarrow{\sim}A\end{aligned}\right.
\]

% slanted
\margin{\deleted{\unsure{et $I$ nilpotent}}}
\uline{Proposition 2} \quad $A,I,B=A'$ comme dessus \inserted{tels que \textcircled{$+$} $H^0(B/A)\xleftarrow{\sim}A$, $H^1(B/A)\simeq0$ (cf \uline{cor. 3})} \inserted{et $I$ \uuline{nilpotent}}, $E'$ un $A'$-module \uline{plat}, muni d'une donnée de descente $D$ relative à $A'/A$, [d'où $E''$, $E'''\ldots$ formant un complexe $\unsure{K|C}_D^*(A'/A;E'\deleted{\unsure{D}})$].

\diagram{\begin{tikzcd}[ampersand replacement=\&,column sep=.27em,row sep=1.2em]
\& \& E'\arrow[rr,equal]\arrow[dd,no head]\arrow[dr,no head] \& \& E''\arrow[rr,shift left=1,no head]\arrow[rr,no head]\arrow[rr,shift right=1,no head]\arrow[dd,no head]\arrow[dr,no head] \& \& E'''\arrow[dd,no head]\arrow[dr,no head]\& \\
\& E_0\arrow[rr,dotted,no head]\arrow[dd,dashed,no head] \& \& E'_0\arrow[rr,equal]\arrow[dd,no head]\& \& E''_0\arrow[rr,shift left=1,no head]\arrow[rr,no head]\arrow[rr,shift right=1,no head]\arrow[dd,no head]\& \& E'''_0\arrow[dd,no head]\\
A\arrow[rr,no head]\arrow[dr,no head]\& \& A'\arrow[rr,equal]\arrow[dr,no head]\& \& A''\arrow[rr,shift left=1,no head]\arrow[rr,no head]\arrow[rr,shift right=1,no head]\arrow[dr,no head]\& \& A'''\arrow[dr,no head]\& \\
\& A_0\arrow[rr,no head]\& \& A'_0\arrow[rr,equal]\& \& A''_0\arrow[rr,shift left=1,no head]\arrow[rr,no head]\arrow[rr,shift right=1,no head]\& \& A'''_0
\end{tikzcd}}

Supposons que la donnée de descente sur $E'_0$ relative à $A'_0/A_0$ déduite par changement de base $A_0\longrightarrow A'_0$ soit effective [pour la catégorie fibrée des Modules plats], i.e. $\exists E_0$ module \uline{plat} sur $A_0$, et un isom compatible aux descentes $E_0\otimes_{A_0}A'_0\simeq E'_0$. Alors on a effectivité

\textcircled{$+$} Ces conditions \deleted{\unsure{impliquent}} signifient que $\operatorname{Spec}B\longrightarrow\operatorname{Spec}A$ est un morphisme de descente pour la cat. fib. des modules plats, et équivalent : \unsure{le $=$ avec|la $=$ cond. pour} $B_0/A_0$ par \uline{cor. 3}
% archive sheet number: 4
\scanpage{6}
également pour la donnée de descente sur $E'$ relt $A'/A$, i.e. si $E\subset E'$ est le sous-$A$-module des \uline{invariants} de $E'$, $E=H^0_D(A'/A;E')$, alors $E$ est \uline{plat} sur $A$ et $\deleted{\unsure{H^0}}E\otimes_A A'\longrightarrow E'$ est un isomorphisme.

\uline{Dém}. \inserted{\uline{Supposons d'abord} $I^2=0$} On voit d'abord qu'on a une suite \deleted{\unsure{courte de complexes|longue de complexes}} exacte de $A$-modules simpliciaux (sans hypothèse d'effectivité pour $E'_0$ rel à $A'_0/A'$)
\[
0\longrightarrow IC_D^*(A'/A,E')\longrightarrow C_D^*(A'/A,E')\longrightarrow C_D^*(A'_0/A_0,E'_0)\longrightarrow0
\]
\deleted{$C_D^*(A'\unsure{,|/}$} et on a pour tt $n$, en raison de la \uline{platitude} de $E'/A'$ donc $E^{(n)}/A^{(n)}$, et de $I^2=0$
\[
IC_D^n(A'/A,E')\simeq \underline{E_0^{(n)\deleted{\unsure{+1|*}}}}\otimes_{A_0^{(n)}}\underset{\substack{\Vert\\ I}}{\underline{IA^{(n)}}}\sim
\tag*{(*)}
\]
[où $A^{(n)}=\overset{n+1}{\bigotimes}_A A'$, $E^{(n)}=E\otimes_A A^{(n)}$ etc]. Grâce à l'effectivité pour $E'_0$ rel $A'_0/A_0$, on peut écrire $(*)$ sous la forme
\[
(\underline{E_0\otimes_{A_0} A_0^{(n)}})\otimes_{A_0^{(n)}}IA^{(n)}
\unsure{=|\simeq}E_0\otimes_{A_0}(\underset{\substack{\Vert\\I}}{\underline{IA^{(n)}}})\simeq\underline{E_0\otimes_{A_0}I}
\]
d'où la \deleted{\unsure{conclusion}} suite exacte de $A$-modules simpliciaux :
\[
0\longrightarrow \underline{\underline{C^*_{E_0\otimes_{A_0}I}}}\longrightarrow C_D^*(A'/A,E')\longrightarrow C_D^*(A'_0/A_0,E'_0)\longrightarrow0
\]
% archive sheet number: 5
\scanpage{7}
On en conclut la suite exacte
\[
0\longrightarrow E_0\otimes_{A_0}I\longrightarrow
\underset{\substack{\Vert\\E}}{H_D^0(A'/A,E')}\longrightarrow
\underset{\substack{\Vert\\E_0}}{H_D^0(A'_0/A_0,E'_0)}\longrightarrow0
\]
compte tenu que
\[
H^{\deleted{\unsure{0}}i}(C_{E_0\otimes_{A_0}I})=
\left\{\begin{aligned}
&E_0\otimes_{A_0}I &&\text{pour }i=0\\
&0 &&\text{pour }i=1\\
&&&(\text{et }=\ i\ne0)
\end{aligned}\right.
\]
Cela prouve d'abord que $E\longrightarrow E_0$ est \uuline{surjectif}, \deleted{donc $\operatorname{Im}(E\longrightarrow E')$ engendre le $A'$-module $E'$, puis} et de façon précise que

a) $\operatorname{Ker}(E\mathrel{\unsure{-|\longrightarrow}} E_0)=IE$, donc $E\longrightarrow E_0$ induit un \uline{isomorphisme} $E/IE=E\otimes_A A_0\xrightarrow{\sim}E_0$

b) $\underline{E/IE}\otimes_{A_0}I\longrightarrow IE$ un \uline{isomorphisme}, donc grâce à a) et $E_0$ plat sur $A_0$, $E$ est \uuline{plat} sur $A$

Considérons maintenant \deleted{le diagramme \unsure{comm.}} \inserted{l'hom}
\diagram{\begin{tikzcd}[ampersand replacement=\&]
E\otimes_A A'\arrow[r]\arrow[d,"\deleted{\mathrm{surj}}","\diagdown"{description,fill=none,inner sep=0pt}] \& E'\arrow[d,"\diagdown"{description,fill=none,inner sep=0pt}]\\
\deleted{E_0\otimes_{A_0}A'_0}\arrow[r,"\deleted{\sim}","\diagdown"{description,fill=none,inner sep=0pt}] \& \deleted{E'_0}
\end{tikzcd}}
c'est un hom de $A'$-modules \uline{plats} qui \deleted{on lit sur \unsure{le diagramme|les diagrammes} que le composé $E\otimes_A A'\longrightarrow E'\longrightarrow E'_0$ est surjectif, donc par Nakayama sur $A'$ (muni de $IA'$ nilpotent) $E\otimes_A A'$} \unsure{devient|donne} un isom. par réduction mod $IA'$, donc est un isomorphisme.
% archive sheet number: 6
\scanpage{8}
Dans le cas général
\[
I^n=0
\]
on procède par récurrence sur $n$, en notant que c'est trivial pour $n=\unsure{1|0}$, déjà fait pour $n=2$, on peut supposer $n>2$ et la chose \deleted{faite} pour $n'<n$. \deleted{\unsure{Cela|Il}} Il suffit alors d'utiliser l'hyp de récurrence et le

\uline{Lemme} \quad Si $I,J$ sont deux idéaux de $A$ tels que $\lambda\mapsto\lambda 1_B$ est une bijection de $I$ ($J$) sur un idéal de $B$, alors $IJ$ a la $=$ propriété, donc aussi les puissances \inserted{$I^n$} de $I$, \unsure{avec} $n\geq1$. \deleted{Si de plus}

$+$ \quad (lemme trivial)

\uline{Corollaire} 1 \quad Soient $A,I,B\deleted{=A'}$ comme dessus [$A\subset B$, $I\subset A$ un idéal \uline{nilpotent} de $A$ tel que $IB=\unsure{B|I}$] \deleted{et tel que} supposons $B_0/A_0$ \deleted{fid. plat} \inserted{morphisme de descente effective pour la cat. \deleted{\unsure{fib}}fibrée des modules plats, \uline{p.ex.} $B_0/A_0$ fid plat}. Alors $\operatorname{Spec}B\longrightarrow\operatorname{Spec}A$ est un morphisme de descente \uline{effective} pour la catégorie fibrée des modules \uline{plats}.
% archive sheet number: 7
\scanpage{9}
On peut $=$ noter le corollaire immédiat

\uline{Corollaire} 2 \quad Soient $\deleted{\unsure{A|I|X}} X$ un préschéma, $I$ \deleted{\unsure{un}} Idéal quasi-cohérent \uline{nilpotent} sur $X$, $X'=\operatorname{Spec}(A')$ un $X$-préschéma, affine sur $X$, tel que $I\longrightarrow A'$ soit un isomorphisme de $I$ sur un \uline{Idéal} de $A'$, et tel que $A'/IA'$ soit fid plat sur $\underline{\mathcal O}_X/I$ [i.e. $\operatorname{Spec}(A'/IA')=X'_0$ fid plat sur $V(I)=X_0$]. Alors $X'\longrightarrow X$ est un \uline{morphisme de descente} \uline{effective} pour la catégorie \inserted{fibrée} des Modules \inserted{quasi coh.} plats.

On \deleted{\unsure{peut en dire}} \inserted{en} déduit

\uline{Corollaire} 3 \quad \deleted{Soit $A$ \illegible{2}} \inserted{Sous les conditions précédentes}, pour tt morphisme \inserted{plat} de préschémas $\deleted{\unsure{Y\longleftarrow X|Y\longrightarrow X}}$ \inserted{$X\longleftarrow Y$} \unsure{posant|formant} $\deleted{\unsure{Z}}Y'=X'\times_X Y$, \deleted{\unsure{on}} la projection $Y'\mathrel{\unsure{-|\longrightarrow}}Y$ est un morphisme de descente \unsure{eff.} universelle pour la catégorie fibrée des Modules quasi coh. plats.
% archive sheet number: 8
\scanpage{10}
Ce résultat s'applique en particulier lorsque $X=\operatorname{Spec}A$, $X'=\operatorname{Spec}A'$, où $A,A'$ satisfont les conditions du corollaire 1 (avec $A'_0/A_0$ fid plat), quand $Y$ est un $A$-préschéma quelconque, et $Y'=Y\otimes_A A'$.

\figure{\begin{tikzpicture}\draw (0,0)--(3,0); \draw (1.02,-.11)--(2.15,-.11);\end{tikzpicture}}

\deleted{\uline{Proposition 3}} \unsure{Revenons|Reprenons} au diagramme commutatif du Cor. 1 de prop. 1, en \deleted{\illegible{1} \unsure{supposant}} prenant $F$ réduit à un pt, on trouve un diagramme commutatif
\diagram{\begin{tikzcd}[ampersand replacement=\&,column sep=small]
0\arrow[r,no head]\& I\arrow[r]\arrow[d,Rightarrow]\& \overset n{\bigotimes}_A B\arrow[r,"\unsure{\varphi_n|\phi_n|\mathrm{pr}}"]\arrow[d]\&\overset n{\bigotimes}_{A_0}B_0\arrow[r,no head]\arrow[d]\&0\\
0\arrow[r,no head]\&I\arrow[r]\&B\arrow[r]\&B_0\arrow[r,no head]\&0
\end{tikzcd}}
on trouve \deleted{le diagramme commutatif} \inserted{\unsure{Vu|par} le \uline{diagramme du serpent}}

\uline{Proposition} 3 \quad \inserted{$A,I,B$ comme ds Prop. 1.} Soient : $\Delta_{B/A}^{(n)}=\operatorname{Ker}(\overset n{\bigotimes}_A B\mathrel{\unsure{-|\longrightarrow}} B)$, \deleted{$0\mathrel{-}I\mathrel{-}\unsure{0|U}$} et définissons de $=$ $\Delta_{B_0/A_0}^{(n)}$, alors l'hom $\overset n{\bigotimes}_A B\longrightarrow\overset n{\bigotimes}_{A_0}B_0$ induit un \uline{isom}
\[
\Delta_{B/A}^{(n)}\xrightarrow{\sim}\Delta_{B_0/A_0}^{(n)}
\]
% archive sheet number: 9
\scanpage{11}
\uline{Corollaire} 1 \quad Il y a correspondance biunivoque entre les idéaux $J$ de $\overset n{\bigotimes}_A B$ contenus dans l'idéal d'augmentation, et les idéaux $J^{\unsure{0|o}}$ de $\overset n{\bigotimes}_{A_0}B_0$ contenus dans l'\unsure{idéale|idéal} d'augmentation, donnée par
\[
\left\{\begin{aligned}
J&\longmapsto\unsure{\rho_n|p_n}(J)\\
\rho_n^{-1}(J^0)\cap\Delta_{B/A}^{(n)}&\longleftarrow J^0
\end{aligned}\right.
\]
Pour deux idéaux \inserted{$J,J^0$} qui se correspondent, l'hom naturel $J\longrightarrow J^{\unsure{0|2}}$ est bijectif.

Noter que la correspondance précédente est compatible avec les relations d'ordre, \deleted{et \unsure{la}} \inserted{\unsure{pour|puis}} avec la formation de $+$ et $-$. D'autre part, les \deleted{\unsure{formations de ces}} \inserted{\unsure{applications|bijections|identifications} précédentes} (pour $n$ variable) commutent aux opérations simpliciales, dans un sens évident. \inserted{(cf Prop 1, cor 1)}

On conclut de \deleted{\unsure{ceci}} \inserted{ces remarques}

\uline{Corollaire} 2 \quad \unsure{(a)|($\alpha$)} Il y a correspondance biunivoque (induite par celle du cor. 1 pour $n=2$) entre relations d'équivalence dans
% archive sheet number: 10
\scanpage{12}
$\operatorname{Spec} B$ sur $\operatorname{Spec} A$, et celles de $\operatorname{Spec} B_0$ sur $\operatorname{Spec} A_0$. \textcircled{b} Cette correspondance est \unsure{compatibles|compatible} avec les relations d'ordre, \deleted{\unsure{et en particulier}} \inserted{en particulier la relation} d'équivalence la plus (la moins) fine de \deleted{plus fine. On en conclut :} $\operatorname{Spec} B$ sur $\operatorname{Spec} A$, correspond à la relation d'équiv. la plus (la moins) fine de $\operatorname{Spec} B_0$ sur $\operatorname{Spec} A_0$.

[\,\textcircled{c} Si $J,J^{\unsure{0|o}}$ sont des idéaux d'\unsure{équivalences|équivalence} qui se correspondent, et $\deleted{\unsure{A|\Delta}} C$ ($C^0$) les ss-anneaux de $B$ ($B_0$) des invariants, alors $C^0\simeq C/IC\simeq C/I$, [de façon plus précise, si $\delta:B\longrightarrow D=B\otimes_A B/J$ et $\delta_0:B_0\longrightarrow D_0$ sont les homs \unsure{cobords|obvious}, alors on a
% slanted
\margin{expliciter plus \unsure{-\,-\,-\,-} \deleted{\unsure{équivalence ici} \unsure{de} $\unsure{ID|\delta|D}$}}
$\delta(b)\in ID$ sss $\deleted{b\in IB}\delta(b)=0$, i.e. $\delta_0(b_0)=0$ sss $\delta(b)=0$ i.e $b_0\in C_0$ sss $b\in C$], donc on écrira aussi $C_0$ aussi. Pour que $J$ soit effectif, il f et s. que $J^0$ le soit \ldots\ ]
% archive sheet number: 11
\scanpage{13}
\textcircled{2} Résultats auxiliaires sur les relations d'équivalence

\uline{Proposition 1} \quad Soient $A$ un anneau, $\unsure{m|\mathfrak m}$ un idéal maximal de $A$, $A'$ une $A$-algèbre, contenant $A$, tel que $A'/A$ soit un $A$-module isomorphe à $A/m=k$

Sous ces conditions

(i) Il \deleted{\unsure{n'}}y a dans $\operatorname{Spec}A'$ sur $\operatorname{Spec}A$ exactement deux relations d'\unsure{équivalences|équivalence}, la plus fine \deleted{et} [correspondant à la diagonale de $\operatorname{Spec}A'$ sur $\operatorname{Spec} A$] et la moins fine [correspondant au produit fibré]

(ii) la suite
\[
\deleted{\unsure{0\longrightarrow}}A\longrightarrow A'\rightrightarrows A'\otimes_A A'
\]
est exacte, \deleted{\unsure{plus pr.|plus précis}} i.e. $H^0(A'/A)\xleftarrow{\sim}A$, et $H^i(A'/A)=0$ pour $i\ne0$

(iii) $\operatorname{Spec}A'\longrightarrow\operatorname{Spec}A$ est un morphisme de descente effective pour la catégorie fibrée des Modules \unsure{qu. coh.:|qu. coh.} plats, et reste \unsure{telle|telles|tel} pour \unsure{toute|toutes} \unsure{extension|extensions} de la base $X\longrightarrow\operatorname{Spec}A$).
% archive sheet number: 12
\scanpage{14}
\uline{Démonstration} \quad $A,A'$ et l'idéal $\unsure{m|\mathfrak m}$ satisfont les conditions de prop 1, de plus $A_0^{\unsure{2|\deleted{2}}}=A/m=k$ est un corps, et $A'/mA'=A'/m\simeq A'_0$ est un vectoriel de dim 2 sur $k$, il est en particulier fid plat sur $k$. Alors (ii) et (iii) résultent du N\textsuperscript{o} 1, Prop. 1, cor. 5, et prop. 2 \inserted{cor. 1 et} cor. 3.

Pour (i), utilisant N\textsuperscript{o} 1, Prop 3, cor. 2, on est ramené au cas où $A=k$, donc $A'$ est une algèbre de rang 2 sur $k$. \inserted{On peut supposer $k$ alg. clos.} On doit distinguer deux cas

(i) \inserted{$A'$ réduit, d'où :} $A'\simeq k\times k$, alors la question est purement ensembliste ; il y a exactement 2 relations d'équivalence sur un ens à 2 éléments.

(ii) $A'\deleted{\simeq k}$ non réduit, alors $A'\simeq k[\varepsilon]$, avec $\varepsilon^2=0$. (\unsure{nbs|nbrs|nbres} duaux).
% archive sheet number: 13
\scanpage{15}
\uline{Proposition} 2 \quad \unsure{Soient|Soit} $A$ un anneau local artinien, $A'$ une $A$-algèbre finie, $A''=(A'\otimes_A A')/J$, où $J$ est un idéal d'équivalence, \deleted{On suppose que} \inserted{$\Delta=\operatorname{Ker}(A''\mathrel{\unsure{-|\longrightarrow}} A')$}. On suppose que \deleted{pour tt $x\in A'$}
\diagram{\begin{tikzcd}[ampersand replacement=\&]
A\arrow[r]\& A'\arrow[r,shift left=.5ex,"p_1"]\arrow[r,shift right=.5ex,"p_2"']\& A''
\end{tikzcd}}
est une suite \uline{exacte}, et que pour tt $x'\in A'\mathbin{-}A$, l'élément
\[
\delta(x')=p_2(x')-p_1(x')
\]
est un \uline{générateur} de l'idéal $\Delta$.
% slanted
\inserted{Supposons de plus les extensions résiduelles de $A'/A$ \uline{triviales}.}
% slanted
\margin{Sans l'hypothèse $\uparrow$, il est vrai que $\unsure{m|\mathfrak m}A'=m$, donc on est ramené à \unsure{étudier|discuter} le cas $A=k$}
Sous ces conditions, $\operatorname{long} A'/A\leq1$, et par suite (cf prop 1, (i)) on a $A''\simeq A'\otimes_A A'$.

Soit $m'=\unsure{\mathfrak r|r}(A')$, distinguons deux cas

1\textsuperscript{o}) $A'$ \uline{n'est pas \deleted{local},} \inserted{radiciel sur $A$} \deleted{Alors $\Delta$ n'est pas} \inserted{i.e. $\operatorname{Ker}(A'\otimes_A A'\longrightarrow A')$} n'est pas nilpotent. \uline{Comme $J$ est nilpotent}, il s'ensuit que $\operatorname{Ker}(A''\longrightarrow A')=\Delta$ n'est pas nilpotent. Mais si $x'\in\sqrt{\unsure{A|A'}}=m'$ \deleted{\unsure{c'est}} alors $\delta(x')$ est nilpotent, donc $\delta(x')A''$ est nilpotent, d'où $\delta(x')A''\ne\Delta$, d'où $x'\in A$ par hyp, d'où
\[
m'\subset A
\]
\inserted{et comme $\deleted{\unsure{mA'|mA}}m'\cap A\subset m\subset mA'\subset m'$, on a} d'où
\[
m'=m=mA'
\]
% archive sheet number: 14
\scanpage{16}
Donc on est sous les conditions du N\textsuperscript{o} 1, avec $A_0=A/\unsure{m|\mathfrak m}=k$ un corps. \deleted{Comme} Nous savons par N\textsuperscript{o} 1 prop. 3, cor. 2 qu'on aura
% slanted
\margin{dégager un lemme de réduction}
\diagram{\begin{tikzcd}[ampersand replacement=\&]
A \arrow[r] \arrow[d] \& A' \arrow[r,shift left=.5ex] \arrow[r,shift right=.5ex] \arrow[d] \& A'' \arrow[d] \\
A_0 \arrow[r] \& A'_0 \arrow[r,shift left=.5ex] \arrow[r,shift right=.5ex] \& A''_0
\end{tikzcd}}
avec la deuxième ligne exacte ; notons qu'on aura $A'_0=A'/mA'=A'/m'$, $A'_0$ est \uline{réduit}. De plus,
\deleted{\unsure{quand|quant|posant}}
\[
\begin{aligned}
\Delta_0=\operatorname{Ker}(A''_0\mathrel{\unsure{-|\longrightarrow}}A'_0)&\simeq\operatorname{Ker}(A''/mA''\longrightarrow A'/mA'),\\
&\simeq\operatorname{Ker}(A''/mA''\mathrel{-}A''/\Delta+mA'')\\
&\simeq(\Delta+mA'')/mA''\simeq\Delta/\Delta\cap mA'',
\end{aligned}
\]
et notant que $\Delta\simeq\operatorname{Ker}(A''\mathrel{-}A')\longrightarrow\Delta_0\simeq\operatorname{Ker}(A''_0\mathrel{-}A'_0)$ est \uline{surjectif}, on voit que \deleted{\unsure{$A$|$\Delta$}} l'hypothèse faite pour $A'/A$ implique que pour tt $x'_0\in A'_0-A_0$, on a $\delta_0(x'_0)\in A'_0$ est un \uline{générateur} de $\deleted{\unsure{A''}}\Delta_0$. Utilisant le fait que les composants locaux de $A'_0$ \deleted{\unsure{$=$|\text{et}}} [qui sont des corps] sont des extensions \uline{triviales} \inserted{de $k=A_0$}, \deleted{et} donc $A'_0\simeq\underbrace{k\times\ldots\times k}_{n}$, avec $n\geq2$, et prenant $x'_0=(1,0,\ldots,0)$ on trouve $n=2$, O.K.
% archive sheet number: 15
